\chapter{Axioms of the Information--Spinor Ansatz}
\label{chap:axioms}

\section{Purpose of the axiomatization}

The preceding chapters established a network of exact correspondences. To convert that network into a research programme, the assumptions must be stated in a form that can be tested, weakened, or rejected. The axioms below are not asserted as properties already possessed by the zeta function. They define the class of models under consideration.

\section{The binary state bundle}

Let $\Sstrip$ be the open critical strip and let $\mathbb P(\C^2)$ be the projective space of rays in a two-dimensional complex Hilbert space. A representative state is written
\begin{equation}
 \ket{\Psi(s)}
 =e^{i\chi(s)}\left(
 \sqrt{\sigma}\ket{0}
 +e^{i\Phi(s)}\sqrt{1-\sigma}\ket{1}
 \right),
 \qquad s=\sigma+it.
 \label{eq:ansatz-state}
\end{equation}
The overall phase $\chi$ is gauge; the relative phase $\Phi$ affects the detector amplitude.

\begin{axiom}[Normalized complement weights]
\label{ax:weights}
The squared channel norms of the state associated with $s=\sigma+it$ are $\sigma$ and $1-\sigma$.
\end{axiom}

This axiom is the most direct identification of the analytic real part with binary balance. It is natural under $\sigma\leftrightarrow1-\sigma$, but it is not forced by classical zeta theory. A more advanced model may derive these weights from a kernel or operator rather than assume them.

\begin{axiom}[Equal-gain detector]
\label{ax:detector}
There exists a fixed detector ray $\ket{+}=(\ket0+\ket1)/\sqrt2$ whose amplitude is the normalized zero readout:
\begin{equation}
 R(s)=\braket{+}{\Psi(s)}.
 \label{eq:readout}
\end{equation}
\end{axiom}

The equality of detector gains is essential. As shown in \cref{chap:cancellation,app:asymmetric}, an asymmetric detector selects a line other than $1/2$.

\begin{postulate}[Zero equivalence]
\label{post:zero-equivalence}
There is a non-vanishing scalar function $g(s)$ on $\Sstrip$ such that
\begin{equation}
 \xi(s)=g(s)R(s),
 \label{eq:zero-factorization}
\end{equation}
or, in a weaker version sufficient for the location theorem,
\begin{equation}
 \xi(s)=0\quad\Longleftrightarrow\quad R(s)=0.
 \label{eq:zero-equivalence}
\end{equation}
Multiplicity preservation is required for a full analytic equivalence.
\end{postulate}

This postulate is SOH-C001 in explicit form. It is the unresolved bridge.

\begin{axiom}[Involution covariance]
\label{ax:covariance}
The projective state map is compatible with the critical-strip involution:
\begin{equation}
 [\Psi(\J s)]=[X\overline{\Psi(s)}],
 \label{eq:ansatz-covariance}
\end{equation}
where $X$ swaps the channel basis.
\end{axiom}

Covariance exchanges the weights $\sigma$ and $1-\sigma$. On the critical line, where $s=\J s$, it enforces equal component norms up to projective gauge.

\begin{axiom}[Canonical construction]
\label{ax:canonical}
The map $s\mapsto[\Psi(s)]$, the detector, and the factor $g$ are determined from classical analytic data such as $\xi$, the theta kernel, an explicit Hilbert-space transform, or a specified operator. They may not use tabulated zero locations or a rule conditioned on $\xi(s)=0$.
\end{axiom}

\section{Minimal and strong versions}

The \emph{minimal ansatz} consists of \cref{ax:weights,ax:detector} and the zero equivalence \cref{eq:zero-equivalence}. It is sufficient for a conditional proof of RH. The \emph{strong ansatz} adds covariance, regularity, multiplicity preservation, and an explicit construction.

The distinction matters because the minimal ansatz can be satisfied by artificial maps. For example, one can choose a phase from the value of $\xi(s)$ so that the readout vanishes whenever $\xi(s)$ does. Such a map restates the zero set without explaining it. The strong ansatz is intended to exclude this.

\section{Gauge-invariant formulation}

Let $P_+=\ket{+}\bra{+}$ be the detector projector. The zero readout condition is equivalent to
\begin{equation}
 \bra{\Psi(s)}P_+\ket{\Psi(s)}=0.
 \label{eq:projector-zero}
\end{equation}
Because $P_+$ is positive semidefinite, this occurs exactly when $P_+\ket{\Psi(s)}=0$. The expression is gauge invariant and non-negative. This suggests a possible positivity route: construct a zeta-dependent state for which a positive expectation value is proportional to $|\xi(s)|^2$ up to a non-vanishing factor.

Specifically, a strong realization could seek
\begin{equation}
 |\xi(s)|^2=w(s)\bra{\Psi(s)}P_+\ket{\Psi(s)},
 \qquad w(s)>0\quad(s\in\Sstrip).
 \label{eq:positive-factorization}
\end{equation}
This avoids phase gauge issues and preserves the zero set. It is non-holomorphic, as any modulus-square identity must be, but it may arise naturally from a reproducing kernel or transfer operator.

\section{Covariance constraints on the phase}

Using the representative \cref{eq:ansatz-state}, the swap-conjugation operation gives
\[
 X\overline{\ket{\Psi(s)}}
 \sim \sqrt{1-\sigma}\ket0+e^{i\Phi(s)}\sqrt\sigma\ket1,
\]
up to an overall phase. At $\J s=(1-\sigma)+it$, this matches the required exchanged weights. Covariance therefore constrains the relative phase at paired points but does not necessarily force
\[
 \Phi(\J s)=\Phi(s).
\]
The precise relation depends on gauge conventions. A projective formulation is safer than equating raw phases.

On the fixed line, covariance requires the ray to be invariant under swap-conjugation. This restricts it to a one-parameter family of balanced rays. The detector zero selects one member of that family. Thus the symmetry and the zero condition remain complementary pieces.

\section{A hierarchy of construction targets}

A candidate bridge can be evaluated against the following hierarchy:
\begin{enumerate}
  \item \textbf{Zero-set map:} correct zeros, perhaps engineered.
  \item \textbf{Covariant map:} respects $\J$ and conjugation.
  \item \textbf{Metric map:} derives the channel weights and equal detector gain.
  \item \textbf{Analytic map:} preserves multiplicity and varies with controlled regularity.
  \item \textbf{Operator map:} arises from a self-adjoint or positive object.
  \item \textbf{Arithmetic map:} retains Euler-product or explicit-formula information.
\end{enumerate}
A map achieving only the first level has little explanatory value. A map reaching the fifth level would be close to a genuine RH mechanism.

\section{The role of the imaginary coordinate}

The axioms determine only the horizontal balance coordinate. The imaginary part $t$ must enter through phase, dynamics, or spectral quantization. In a state model one expects
\[
 \Phi(s)=\Phi(\sigma,t),
\]
with zeros occurring when the phase reaches opposition simultaneously with equal weights. On the critical line, equal weights are automatic, so the zero ordinates are the values of $t$ for which the phase or detector condition quantizes.

This suggests the interpretation
\[
 \sigma\quad\text{as balance},
 \qquad
 t\quad\text{as spectral phase or frequency}.
\]
It is a useful decomposition, not an established physical ontology.

\status{The axioms define the ansatz. Only the elementary consequences are proved. The zero-equivalence and canonical-construction requirements remain unfulfilled.}
